% Cover design and front matter.
\definecolor{covercream}{HTML}{FFF2A0}
\definecolor{coverpaper}{HTML}{FFFEF4}
\definecolor{coveryellow}{HTML}{FFDA48}
\definecolor{coverblue}{HTML}{003568}
\definecolor{coverolive}{HTML}{827032}
\newcommand{\GTMcover}{%
  \thispagestyle{empty}\noindent\hypertarget{p-0}{}\null%
  \begin{tikzpicture}[remember picture,overlay,x=1bp,y=1bp]
  \begin{scope}[shift={(current page.south west)}]
  \fill[coveryellow] (0,0) rectangle (439.37,666.14);
  \shade[left color=covercream,right color=coverpaper] (0,452) rectangle (439.37,666.14);
  \fill[covercream,opacity=.72] (0,0) rectangle (15.5,452);
  \draw[white,line width=.7bp] (15.5,0) -- (15.5,666.14);
  \draw[white,line width=.7bp] (0,452) -- (439.37,452);
  \draw[white,line width=.7bp] (0,357) -- (15.5,357);
  \node[anchor=north west,text=coverolive,inner sep=0,font=\sffamily\bfseries\fontsize{12}{14}\selectfont] at (38,642) {Graduate Texts in Mathematics};
  \node[anchor=north east,text=coveryellow,inner sep=0,font=\sffamily\bfseries\fontsize{70}{72}\selectfont] at (411,656) {GTM};
  \node[anchor=south west,text=coverblue,inner sep=0,font=\sffamily\fontsize{19}{22}\selectfont] at (38,466) {Robin Hartshorne};
  \node[anchor=north west,text=coverblue,inner sep=0,font=\sffamily\fontsize{49}{52}\selectfont,align=left] at (37,437) {Algebraic\\Geometry};
  \node[anchor=west,text=coverblue,inner sep=0,font=\sffamily\small] at (38,58) {Unofficial retypeset edition};

  \node[anchor=south east,inner sep=0] at (396,39) {\includegraphics[width=107bp]{assets/springer-mark.pdf}};
  \end{scope}
  \end{tikzpicture}\clearpage}

\GTMcover
% Contents from body/front-matter.tex.
% Front matter transcribed from the original printed pages.
\newcommand{\FrontBlock}[4]{%
  \put(#1,#2){\makebox(0,0)[lt]{\parbox[t]{#3bp}{#4}}}}
\newcommand{\FrontProse}{\normalsize\setlength{\parindent}{11pt}}
\newcommand{\CatalogEntry}[3]{%
  \par\hangindent=18bp\hangafter=1\noindent
  \makebox[18bp][l]{#1}\textsc{#2}. #3\par}
\newcommand{\CatalogStyle}{\fontsize{8}{9.55}\selectfont\raggedright}
\newcommand{\PublisherTitle}{%
  \raisebox{-.28\height}{\includegraphics[trim=0 0 85bp 0,clip,width=18bp]{assets/springer-mark.pdf}}\enspace Springer}

\BeginTypeSheet{1}{}{}{0}
\FrontBlock{53}{593}{333.37}{\raggedleft
  {\large\bfseries Graduate Texts in Mathematics}\qquad {\Huge 52}\par\vspace{17pt}
  \textit{Editorial Board}\par
  S. Axler\qquad F.W. Gehring\qquad K.A. Ribet}
\EndSheet

\BeginTypeSheet{2}{}{}{0}
\FrontBlock{53}{637}{333.37}{\Large\bfseries Graduate Texts in Mathematics}
\FrontBlock{53}{604}{333.37}{\raggedleft\small\textit{most recent titles in the GTM series}}
\FrontBlock{53}{582}{165}{\CatalogStyle
\CatalogEntry{129}{Fulton/Harris}{Representation Theory: A First Course.\par\hspace{18bp}\textit{Readings in Mathematics}}
\CatalogEntry{130}{Dodson/Poston}{Tensor Geometry.}
\CatalogEntry{131}{Lam}{A First Course in Noncommutative Rings.}
\CatalogEntry{132}{Beardon}{Iteration of Rational Functions.}
\CatalogEntry{133}{Harris}{Algebraic Geometry: A First Course.}
\CatalogEntry{134}{Roman}{Coding and Information Theory.}
\CatalogEntry{135}{Roman}{Advanced Linear Algebra.}
\CatalogEntry{136}{Adkins/Weintraub}{Algebra: An Approach via Module Theory.}
\CatalogEntry{137}{Axler/Bourdon/Ramey}{Harmonic Function Theory.}
\CatalogEntry{138}{Cohen}{A Course in Computational Algebraic Number Theory.}
\CatalogEntry{139}{Bredon}{Topology and Geometry.}
\CatalogEntry{140}{Aubin}{Optima and Equilibria. An Introduction to Nonlinear Analysis.}
\CatalogEntry{141}{Becker/Weispfenning/Kredel}{Gr\"obner Bases. A Computational Approach to Commutative Algebra.}
\CatalogEntry{142}{Lang}{Real and Functional Analysis. 3rd ed.}
\CatalogEntry{143}{Doob}{Measure Theory.}
\CatalogEntry{144}{Dennis/Farb}{Noncommutative Algebra.}
\CatalogEntry{145}{Vick}{Homology Theory. An Abstract Introduction to Algebraic Topology. 2nd ed.}
\CatalogEntry{146}{Bridges}{Computability: A Mathematical Sketchbook.}
\CatalogEntry{147}{Rosenberg}{Algebraic K-Theory and Its Applications.}
\CatalogEntry{148}{Rotman}{An Introduction to the Theory of Groups. 4th ed.}
\CatalogEntry{149}{Ratcliffe}{Foundations of Hyperbolic Manifolds.}
\CatalogEntry{150}{Eisenbud}{Commutative Algebra with a View Toward Algebraic Geometry.}
\CatalogEntry{151}{Silverman}{Advanced Topics in the Arithmetic of Elliptic Curves.}
\CatalogEntry{152}{Ziegler}{Lectures on Polytopes.}
\CatalogEntry{153}{Fulton}{Algebraic Topology: A First Course.}
\CatalogEntry{154}{Brown/Pearcy}{An Introduction to Analysis.}
\CatalogEntry{155}{Kassel}{Quantum Groups.}
\CatalogEntry{156}{Kechris}{Classical Descriptive Set Theory.}
\CatalogEntry{157}{Malliavin}{Integration and Probability.}
\CatalogEntry{158}{Roman}{Field Theory.}
}
\FrontBlock{225}{582}{162}{\CatalogStyle
\CatalogEntry{159}{Conway}{Functions of One Complex Variable II.}
\CatalogEntry{160}{Lang}{Differential and Riemannian Manifolds.}
\CatalogEntry{161}{Borwein/Erd\'elyi}{Polynomials and Polynomial Inequalities.}
\CatalogEntry{162}{Alperin/Bell}{Groups and Representations.}
\CatalogEntry{163}{Dixon/Mortimer}{Permutation Groups.}
\CatalogEntry{164}{Nathanson}{Additive Number Theory: The Classical Bases.}
\CatalogEntry{165}{Nathanson}{Additive Number Theory: Inverse Problems and the Geometry of Sumsets.}
\CatalogEntry{166}{Sharpe}{Differential Geometry: Cartan's Generalization of Klein's Erlangen Program.}
\CatalogEntry{167}{Morandi}{Field and Galois Theory.}
\CatalogEntry{168}{Ewald}{Combinatorial Convexity and Algebraic Geometry.}
\CatalogEntry{169}{Bhatia}{Matrix Analysis.}
\CatalogEntry{170}{Bredon}{Sheaf Theory. 2nd ed.}
\CatalogEntry{171}{Petersen}{Riemannian Geometry.}
\CatalogEntry{172}{Remmert}{Classical Topics in Complex Function Theory.}
\CatalogEntry{173}{Diestel}{Graph Theory.}
\CatalogEntry{174}{Bridges}{Foundations of Real and Abstract Analysis.}
\CatalogEntry{175}{Lickorish}{An Introduction to Knot Theory.}
\CatalogEntry{176}{Lee}{Riemannian Manifolds.}
\CatalogEntry{177}{Newman}{Analytic Number Theory.}
\CatalogEntry{178}{Clarke/Ledyaev/Stern/Wolenski}{Nonsmooth Analysis and Control Theory.}
\CatalogEntry{179}{Douglas}{Banach Algebra Techniques in Operator Theory. 2nd ed.}
\CatalogEntry{180}{Srivastava}{A Course on Borel Sets.}
\CatalogEntry{181}{Kress}{Numerical Analysis.}
\CatalogEntry{182}{Walter}{Ordinary Differential Equations.}
\CatalogEntry{183}{Megginson}{An Introduction to Banach Space Theory.}
\CatalogEntry{184}{Bollob\'as}{Modern Graph Theory.}
\CatalogEntry{185}{Cox/Little/O'Shea}{Using Algebraic Geometry.}
\CatalogEntry{186}{Ramakrishnan/Valenza}{Fourier Analysis on Number Fields.}
\CatalogEntry{187}{Harris/Morrison}{Moduli of Curves.}
\CatalogEntry{188}{Goldblatt}{Lectures on the Hyperreals: An Introduction to Nonstandard Analysis.}
\CatalogEntry{189}{Lam}{Lectures on Modules and Rings.}
\CatalogEntry{190}{Esmonde/Murty}{Problems in Algebraic Number Theory.}
}
\EndSheet

\BeginTypeSheet{3}{}{}{0}
\FrontBlock{53}{596}{333.37}{\Large Robin Hartshorne}
\FrontBlock{53}{540}{333.37}{\Huge\bfseries Algebraic Geometry}
\FrontBlock{53}{77}{333.37}{\Large\PublisherTitle}
\EndSheet

\BeginTypeSheet{4}{}{}{0}
\FrontBlock{53}{563}{333.37}{%
  Robin Hartshorne\\
  Department of Mathematics\\
  University of California\\
  Berkeley, California 94720\\
  USA\par\vspace{28pt}
  \textit{Editorial Board}\par\vspace{5pt}
  \begin{minipage}[t]{163bp}
  S. Axler\\
  Mathematics Department\\
  San Francisco State University\\
  San Francisco, CA 94132\\
  USA
  \end{minipage}\hfill
  \begin{minipage}[t]{165bp}
  K.A. Ribet\\
  Department of Mathematics\\
  University of California at Berkeley\\
  Berkeley, CA 94720-3840\\
  USA\\
  ribet@math.berkeley.edu
  \end{minipage}}
\FrontBlock{53}{347}{333.37}{\fontsize{8}{9.5}\selectfont
  Mathematics Subject Classification (2000): 13-xx, 14A10, 14A15, 14Fxx, 14Hxx, 14Jxx\par\vspace{8pt}
  Library of Congress Cataloging-in-Publication Data\par
  Hartshorne, Robin.\par
  \quad Algebraic geometry.\par\vspace{2pt}
  \quad (Graduate texts in mathematics: 52)\par
  \quad Bibliography: p.\par
  \quad Includes index.\par
  \quad 1. Geometry, Algebraic.\enspace I. Title\enspace II. Series.\par
  QA564.H25\qquad 516'.35\qquad 77-1177\par\vspace{7pt}
  ISBN 978-1-4419-2807-8\quad ISBN 978-1-4757-3849-0 (eBook)\hfill Printed on acid-free paper.\par
  DOI 10.1007/978-1-4757-3849-0\par\vspace{8pt}
  \copyright{} 1977 Springer Science+Business Media, Inc.\par
  Softcover reprint of the hardcover 1st edition 1977\par
  All rights reserved. This work may not be translated or copied in whole or in part without the written permission of the publisher (Springer Science+Business Media, Inc., 233 Spring Street, New York, NY 10013, USA), except for brief excerpts in connection with reviews or scholarly analysis. Use in connection with any form of information storage and retrieval, electronic adaptation, computer software, or by similar or dissimilar methodology now known or hereafter developed is forbidden.\par
  The use in this publication of trade names, trademarks, service marks, and similar terms, even if they are not identified as such, is not to be taken as an expression of opinion as to whether or not they are subject to proprietary rights.\par\vspace{8pt}
  \hspace*{171bp}(ASC/SBA)\par\vspace{8pt}
  15\enspace 14\par\vspace{8pt}
  springeronline.com}
\EndSheet

\BeginTypeSheet{5}{}{}{0}
\FrontBlock{53}{398}{333.37}{\centering\textit{For Edie, Jonathan, and Benjamin}\par}
\EndSheet

\BookMark{0}{Preface}{p-6}
\BeginTypeSheet{6}{vii}{}{1}
\Opening{53}{577}{p-6}{}{Preface}
\FrontBlock{53}{411}{333.37}{\FrontProse
\noindent This book provides an introduction to abstract algebraic geometry using the methods of schemes and cohomology. The main objects of study are algebraic varieties in an affine or projective space over an algebraically closed field; these are introduced in Chapter \Goto{p-15}{I}, to establish a number of basic concepts and examples. Then the methods of schemes and cohomology are developed in Chapters \Goto{p-74}{II} and \Goto{p-215}{III}, with emphasis on applications rather than excessive generality. The last two chapters of the book (\Goto{p-307}{IV} and \Goto{p-370}{V}) use these methods to study topics in the classical theory of algebraic curves and surfaces.

The prerequisites for this approach to algebraic geometry are results from commutative algebra, which are stated as needed, and some elementary topology. No complex analysis or differential geometry is necessary. There are more than four hundred exercises throughout the book, offering specific examples as well as more specialized topics not treated in the main text. Three appendices present brief accounts of some areas of current research.

This book can be used as a textbook for an introductory course in algebraic geometry, following a basic graduate course in algebra. I recently taught this material in a five-quarter sequence at Berkeley, with roughly one chapter per quarter. Or one can use Chapter \Goto{p-15}{I} alone for a short course. A third possibility worth considering is to study Chapter \Goto{p-15}{I}, and then proceed directly to Chapter \Goto{p-307}{IV}, picking up only a few definitions from Chapters \Goto{p-74}{II} and \Goto{p-215}{III}, and \emph{assuming} the statement of the Riemann--Roch theorem for curves. This leads to interesting material quickly, and may provide better motivation for tackling Chapters \Goto{p-74}{II} and \Goto{p-215}{III} later.

The material covered in this book should provide adequate preparation for reading more advanced works such as Grothendieck \Goto{bib-EGA}{[EGA]}, \Goto{bib-SGA}{[SGA]}, Hartshorne \Goto{legacy-b-h5}{[5]}, Mumford [2], [5], or Shafarevich [1].}
\EndSheet

\BeginTypeSheet{7}{viii}{\Goto{p-6}{Preface}}{0}
\FrontBlock{53}{585}{333.37}{\FrontProse
\textit{Acknowledgements}\par
In writing this book, I have attempted to present what is essential for a basic course in algebraic geometry. I wanted to make accessible to the nonspecialist an area of mathematics whose results up to now have been widely scattered, and linked only by unpublished \textquotedblleft{}folklore.\textquotedblright{} While I have reorganized the material and rewritten proofs, the book is mostly a synthesis of what I have learned from my teachers, my colleagues, and my students. They have helped in ways too numerous to recount. I owe especial thanks to Oscar Zariski, J.-P. Serre, David Mumford, and Arthur Ogus for their support and encouragement.

Aside from the \textquotedblleft{}classical\textquotedblright{} material, whose origins need a historian to trace, my greatest intellectual debt is to A. Grothendieck, whose treatise \Goto{bib-EGA}{[EGA]} is the authoritative reference for schemes and cohomology. His results appear without specific attribution throughout Chapters \Goto{p-74}{II} and \Goto{p-215}{III}. Otherwise I have tried to acknowledge sources whenever I was aware of them.

In the course of writing this book, I have circulated preliminary versions of the manuscript to many people, and have received valuable comments from them. To all of these people my thanks, and in particular to J.-P. Serre, H. Matsumura, and Joe Lipman for their careful reading and detailed suggestions.

I have taught courses at Harvard and Berkeley based on this material, and I thank my students for their attention and their stimulating questions.

I thank Richard Bassein, who combined his talents as mathematician and artist to produce the illustrations for this book.

A few words cannot adequately express the thanks I owe to my wife, Edie Churchill Hartshorne. While I was engrossed in writing, she created a warm home for me and our sons Jonathan and Benjamin, and through her constant support and friendship provided an enriched human context for my life.

For financial support during the preparation of this book, I thank the Research Institute for Mathematical Sciences of Kyoto University, the National Science Foundation, and the University of California at Berkeley.

\vspace{16pt}
\noindent August 29, 1977\par
\noindent Berkeley, California\hfill \textsc{Robin Hartshorne}\par}
\EndSheet
% Contents from body/page-008.tex.
% Original PDF page 9, printed page ix
\BookMark{0}{Contents}{p-8}
\ContentsSheet{8}{ix}
\ContentsRow{0}{p-11}{Introduction}{xiii}
\ContentsRow{0}{p-15}{I.  Varieties}{1}
\ContentsRow{1}{sec-I-1}{1\quad Affine Varieties}{1}
\ContentsRow{1}{sec-I-2}{2\quad Projective Varieties}{8}
\ContentsRow{1}{sec-I-3}{3\quad Morphisms}{14}
\ContentsRow{1}{sec-I-4}{4\quad Rational Maps}{24}
\ContentsRow{1}{sec-I-5}{5\quad Nonsingular Varieties}{31}
\ContentsRow{1}{sec-I-6}{6\quad Nonsingular Curves}{39}
\ContentsRow{1}{sec-I-7}{7\quad Intersections in Projective Space}{47}
\ContentsRow{1}{sec-I-8}{8\quad What Is Algebraic Geometry?}{55}
\ContentsRow{0}{p-74}{II.  Schemes}{60}
\ContentsRow{1}{sec-II-1}{1\quad Sheaves}{60}
\ContentsRow{1}{sec-II-2}{2\quad Schemes}{69}
\ContentsRow{1}{sec-II-3}{3\quad First Properties of Schemes}{82}
\ContentsRow{1}{sec-II-4}{4\quad Separated and Proper Morphisms}{95}
\ContentsRow{1}{sec-II-5}{5\quad Sheaves of Modules}{108}
\ContentsRow{1}{sec-II-6}{6\quad Divisors}{129}
\ContentsRow{1}{sec-II-7}{7\quad Projective Morphisms}{149}
\ContentsRow{1}{sec-II-8}{8\quad Differentials}{172}
\ContentsRow{1}{sec-II-9}{9\quad Formal Schemes}{190}
\ContentsRow{0}{p-215}{III.  Cohomology}{201}
\ContentsRow{1}{sec-III-1}{1\quad Derived Functors}{202}
\ContentsRow{1}{sec-III-2}{2\quad Cohomology of Sheaves}{206}
\ContentsRow{1}{sec-III-3}{3\quad Cohomology of a Noetherian Affine Scheme}{213}
\EndContentsSheet
% Contents from body/page-009.tex.
% Original PDF page 10, printed page x
\ContentsSheet{9}{x}
\ContentsRow{1}{sec-III-4}{4\quad \v{C}ech Cohomology}{218}
\ContentsRow{1}{sec-III-5}{5\quad The Cohomology of Projective Space}{225}
\ContentsRow{1}{sec-III-6}{6\quad Ext Groups and Sheaves}{233}
\ContentsRow{1}{sec-III-7}{7\quad The Serre Duality Theorem}{239}
\ContentsRow{1}{sec-III-8}{8\quad Higher Direct Images of Sheaves}{250}
\ContentsRow{1}{sec-III-9}{9\quad Flat Morphisms}{253}
\ContentsRow{1}{sec-III-10}{10\quad Smooth Morphisms}{268}
\ContentsRow{1}{sec-III-11}{11\quad The Theorem on Formal Functions}{276}
\ContentsRow{1}{sec-III-12}{12\quad The Semicontinuity Theorem}{281}
\ContentsRow{0}{p-307}{IV.  Curves}{293}
\ContentsRow{1}{sec-IV-1}{1\quad Riemann\textendash{}Roch Theorem}{294}
\ContentsRow{1}{sec-IV-2}{2\quad Hurwitz\textquoteright{}s Theorem}{299}
\ContentsRow{1}{sec-IV-3}{3\quad Embeddings in Projective Space}{307}
\ContentsRow{1}{sec-IV-4}{4\quad Elliptic Curves}{316}
\ContentsRow{1}{sec-IV-5}{5\quad The Canonical Embedding}{340}
\ContentsRow{1}{sec-IV-6}{6\quad Classification of Curves in \texorpdfstring{\ensuremath{\mathbf P^3}}{P3}}{349}
\ContentsRow{0}{p-370}{V.  Surfaces}{356}
\ContentsRow{1}{sec-V-1}{1\quad Geometry on a Surface}{357}
\ContentsRow{1}{sec-V-2}{2\quad Ruled Surfaces}{369}
\ContentsRow{1}{sec-V-3}{3\quad Monoidal Transformations}{386}
\ContentsRow{1}{sec-V-4}{4\quad The Cubic Surface in \texorpdfstring{\ensuremath{\mathbf P^3}}{P3}}{395}
\ContentsRow{1}{sec-V-5}{5\quad Birational Transformations}{409}
\ContentsRow{1}{sec-V-6}{6\quad Classification of Surfaces}{421}
\ContentsRow{0}{p-438}{Appendix A.  Intersection Theory}{424}
\ContentsRow{1}{sec-A-1}{1\quad Intersection Theory}{425}
\ContentsRow{1}{sec-A-2}{2\quad Properties of the Chow Ring}{428}
\ContentsRow{1}{sec-A-3}{3\quad Chern Classes}{429}
\ContentsRow{1}{sec-A-4}{4\quad The Riemann\textendash{}Roch Theorem}{431}
\ContentsRow{1}{sec-A-5}{5\quad Complements and Generalizations}{434}
\ContentsRow{0}{p-452}{Appendix B.  Transcendental Methods}{438}
\ContentsRow{1}{sec-B-1}{1\quad The Associated Complex Analytic Space}{438}
\ContentsRow{1}{sec-B-2}{2\quad Comparison of the Algebraic and Analytic Categories}{440}
\ContentsRow{1}{sec-B-3}{3\quad When is a Compact Complex Manifold Algebraic?}{441}
\ContentsRow{1}{sec-B-4}{4\quad K\"{a}hler Manifolds}{445}
\ContentsRow{1}{sec-B-5}{5\quad The Exponential Sequence}{446}
\EndContentsSheet
% Contents from body/page-010.tex.
% Original PDF page 11, printed page xi
\ContentsSheet{10}{xi}
\ContentsRow{0}{p-463}{Appendix C.  The Weil Conjectures}{449}
\ContentsRow{1}{sec-C-1}{1\quad The Zeta Function and the Weil Conjectures}{449}
\ContentsRow{1}{sec-C-2}{2\quad History of Work on the Weil Conjectures}{451}
\ContentsRow{1}{sec-C-3}{3\quad The \texorpdfstring{\ensuremath{\ell}}{l}-adic Cohomology}{453}
\ContentsRow{1}{sec-C-4}{4\quad Cohomological Interpretation of the Weil Conjectures}{454}
\ContentsRow{0}{p-473}{Bibliography}{459}
\ContentsRow{0}{p-484}{Results from Algebra}{470}
\ContentsRow{0}{p-486}{Glossary of Notations}{472}
\ContentsRow{0}{p-492}{Index}{478}
\EndContentsSheet
% Contents from body/introduction.tex.
% Introduction transcribed from original printed pages xiii--xvi.
\BookMark{0}{Introduction}{p-11}
\BeginTypeSheet{11}{xiii}{}{1}
\Opening{53}{577}{p-11}{}{Introduction}
\FrontBlock{53}{414}{333.37}{\FrontProse
\noindent The author of an introductory book on algebraic geometry has the difficult task of providing geometrical insight and examples, while at the same time developing the modern technical language of the subject. For in algebraic geometry, a great gap appears to separate the intuitive ideas which form the point of departure from the technical methods used in current research.

The first question is that of language. Algebraic geometry has developed in waves, each with its own language and point of view. The late nineteenth century saw the function-theoretic approach of Riemann, the more geometric approach of Brill and Noether, and the purely algebraic approach of Kronecker, Dedekind, and Weber. The Italian school followed with Castelnuovo, Enriques, and Severi, culminating in the classification of algebraic surfaces. Then came the twentieth-century \textquotedblleft{}American\textquotedblright{} school of Chow, Weil, and Zariski, which gave firm algebraic foundations to the Italian intuition. Most recently, Serre and Grothendieck initiated the French school, which has rewritten the foundations of algebraic geometry in terms of schemes and cohomology, and which has an impressive record of solving old problems with new techniques. Each of these schools has introduced new concepts and methods. In writing an introductory book, is it better to use the older language which is closer to the geometric intuition, or to start at once with the technical language of current research?

The second question is a conceptual one. Modern mathematics tends to obliterate history: each new school rewrites the foundations of its subject in its own language, which makes for fine logic but poor pedagogy. Of what use is it to know the definition of a scheme if one does not realize that a ring of integers in an algebraic number field, an algebraic curve, and a compact Riemann surface are all examples of a \textquotedblleft{}regular scheme of}
\EndSheet

\BeginTypeSheet{12}{xiv}{\Goto{p-11}{Introduction}}{0}
\FrontBlock{53}{598}{333.37}{\FrontProse
\noindent dimension one\textquotedblright{}? How then can the author of an introductory book indicate the inputs to algebraic geometry coming from number theory, commutative algebra, and complex analysis, and also introduce the reader to the main objects of study, which are algebraic varieties in affine or projective space, while at the same time developing the modern language of schemes and cohomology? What choice of topics will convey the meaning of algebraic geometry, and still serve as a firm foundation for further study and research?

\medskip
My own bias is somewhat on the side of classical geometry. I believe that the most important problems in algebraic geometry are those arising from old-fashioned varieties in affine or projective spaces. They provide the geometric intuition which motivates all further developments. In this book, I begin with a chapter on varieties, to establish many examples and basic ideas in their simplest form, uncluttered with technical details. Only after that do I develop systematically the language of schemes, coherent sheaves, and cohomology, in Chapters \Goto{p-74}{II} and \Goto{p-215}{III}. These chapters form the technical heart of the book. In them I attempt to set forth the most important results, but without striving for the utmost generality. Thus, for example, the cohomology theory is developed only for quasi-coherent sheaves on noetherian schemes, since this is simpler and sufficient for most applications; the theorem of \textquotedblleft{}coherence of direct image sheaves\textquotedblright{} is proved only for projective morphisms, and not for arbitrary proper morphisms. For the same reasons I do not include the more abstract notions of representable functors, algebraic spaces, \'etale cohomology, sites, and topoi.

The fourth and fifth chapters treat classical material, namely nonsingular projective curves and surfaces, but they use techniques of schemes and cohomology. I hope these applications will justify the effort needed to absorb all the technical apparatus in the two previous chapters.

As the basic language and logical foundation of algebraic geometry, I have chosen to use commutative algebra. It has the advantage of being precise. Also, by working over a base field of arbitrary characteristic, which is necessary in any case for applications to number theory, one gains new insight into the classical case of base field $\mathbf{C}$. Some years ago, when Zariski began to prepare a volume on algebraic geometry, he had to develop the necessary algebra as he went. The task grew to such proportions that he produced a book on commutative algebra only. Now we are fortunate in having a number of excellent books on commutative algebra: Atiyah--Macdonald \Goto{legacy-b-am}{[1]}, Bourbaki \Goto{legacy-b-bour}{[1]}, Matsumura \Goto{legacy-b-mat}{[2]}, Nagata [7], and Zariski--Samuel \Goto{legacy-b-zs}{[1]}. My policy is to quote purely algebraic results as needed, with references to the literature for proof. A list of the results used appears at the end of the book.

\medskip
Originally I had planned a whole series of appendices---short expository accounts of some current research topics, to form a bridge between the main text of this book and the research literature. Because of limited}
\EndSheet

\BeginTypeSheet{13}{xv}{\Goto{p-11}{Introduction}}{0}
\FrontBlock{53}{608}{333.37}{\FrontProse
\noindent time and space only three survive. I can only express my regret at not including the others, and refer the reader instead to the Arcata volume (Hartshorne, ed. [1]) for a series of articles by experts in their fields, intended for the nonspecialist. Also, for the historical development of algebraic geometry let me refer to Dieudonn\'e [1]. Since there was not space to explore the relation of algebraic geometry to neighboring fields as much as I would have liked, let me refer to the survey article of Cassels [1] for connections with number theory, and to Shafarevich [2, Part III] for connections with complex manifolds and topology.

Because I believe strongly in active learning, there are a great many exercises in this book. Some contain important results not treated in the main text. Others contain specific examples to illustrate general phenomena. I believe that the study of particular examples is inseparable from the development of general theories. The serious student should attempt as many as possible of these exercises, but should not expect to solve them immediately. Many will require a real creative effort to understand. An asterisk denotes a more difficult exercise. Two asterisks denote an unsolved problem.

See \Goto{sec-I-8}{(I, \S8)} for a further introduction to algebraic geometry and this book.

\smallskip
\textit{Terminology}\par
For the most part, the terminology of this book agrees with generally accepted usage, but there are a few exceptions worth noting. A \emph{variety} is always irreducible and is always over an algebraically closed field. In Chapter \Goto{p-15}{I} all varieties are quasi-projective. In \Goto{sec-II-4}{(Ch. II, \S4)} the definition is expanded to include \emph{abstract varieties}, which are integral separated schemes of finite type over an algebraically closed field. The words \emph{curve}, \emph{surface}, and \emph{3-fold} are used to mean varieties of dimension 1, 2, and 3 respectively. But in Chapter \Goto{p-307}{IV}, the word \emph{curve} is used only for a nonsingular projective curve; whereas in Chapter \Goto{p-370}{V} a \emph{curve} is any effective divisor on a nonsingular projective surface. A \emph{surface} in Chapter \Goto{p-370}{V} is always a nonsingular projective surface.

A \emph{scheme} is what used to be called a prescheme in the first edition of \Goto{bib-EGA}{[EGA]}, but is called scheme in the new edition of \Goto{bib-EGA}{[EGA, Ch.~I]}.

The definitions of a \emph{projective morphism} and a \emph{very ample invertible sheaf} in this book are not equivalent to those in \Goto{bib-EGA}{[EGA]}---see (II, \Goto{sec-II-4}{\S4}, \Goto{sec-II-5}{5}). They are technically simpler, but have the disadvantage of not being local on the base.

The word \emph{nonsingular} applies only to varieties; for more general schemes, the words \emph{regular} and \emph{smooth} are used.

\smallskip
\textit{Results from algebra}\par
I assume the reader is familiar with basic results about rings, ideals, modules, noetherian rings, and integral dependence, and is willing to accept or look up other results, belonging properly to commutative algebra}
\EndSheet

\BeginTypeSheet{14}{xvi}{\Goto{p-11}{Introduction}}{0}
\FrontBlock{53}{600}{333.37}{\FrontProse
\noindent or homological algebra, which will be stated as needed, with references to the literature. These results will be marked with an A: e.g., Theorem 3.9A, to distinguish them from results proved in the text.

The basic conventions are these: All rings are commutative with identity element $1$. All homomorphisms of rings take $1$ to $1$. In an integral domain or a field, $0\ne 1$. A \emph{prime ideal} (respectively, \emph{maximal ideal}) is an ideal $\mathfrak{p}$ in a ring $A$ such that the quotient ring $A/\mathfrak{p}$ is an integral domain (respectively, a field). Thus the ring itself is not considered to be a prime ideal or a maximal ideal.

A \emph{multiplicative system} in a ring $A$ is a subset $S$, containing $1$, and closed under multiplication. The \emph{localization} $S^{-1}A$ is defined to be the ring formed by equivalence classes of fractions $a/s$, $a\in A$, $s\in S$, where $a/s$ and $a'/s'$ are said to be \emph{equivalent} if there is an $s''\in S$ such that $s''(s'a-sa')=0$ (see e.g. Atiyah--Macdonald \Goto{legacy-b-am}{[1, Ch.~3]}). Two special cases which are used constantly are the following. If $\mathfrak{p}$ is a prime ideal in $A$, then $S=A-\mathfrak{p}$ is a multiplicative system, and the corresponding localization is denoted by $A_{\mathfrak{p}}$. If $f$ is an element of $A$, then $S=\{1\}\cup\{f^n\mid n\geqslant 1\}$ is a multiplicative system, and the corresponding localization is denoted by $A_f$. (Note for example that if $f$ is nilpotent, then $A_f$ is the zero ring.)

\smallskip
\textit{References}\par
Bibliographical references are given by author, with a number in square brackets to indicate which work, e.g. Serre, [3, p.~75]. Cross references to theorems, propositions, lemmas within the same chapter are given by number in parentheses, e.g. (3.5). Reference to an exercise is given by (Ex.~3.5). References to results in another chapter are preceded by the chapter number, e.g. (II, 3.5), or (II, Ex.~3.5).}
\EndSheet
